\heading{高等数学A（II）-模拟题2-期末}

\noteinfo{题目和答案由AI生成。}

\begin{question}
求解下列初值问题：
\begin{enumerate}[label=（\arabic*）,leftmargin=2.8em,itemsep=0.2em]
\item $xy'+2y=x^3\ln x,\quad x>0,\ y(1)=0$；
\item $x^2y''-3xy'+4y=x^2,\quad x>0,\ y(1)=0,\ y'(1)=1$.
\end{enumerate}
\begin{solution}

\begin{enumerate}[label=（\arabic*）,leftmargin=2.8em,itemsep=0.2em]
\item 方程为一阶线性方程：
\[
y'+\frac{2}{x}y=x^2\ln x.
\]
积分因子为 $x^2$，故
\[
(x^2y)'=x^4\ln x.
\]
积分得
\[
x^2y=\int x^4\ln x\,dx=
\frac{x^5}{5}\ln x-\frac{x^5}{25}+C.
\]
由 $y(1)=0$ 得 $C=1/25$，故
\ansmath{y(x)=\frac{x^3}{5}\ln x-\frac{x^3}{25}+\frac{1}{25x^2}.}

\item 这是 Euler 方程。令 $x=e^t$，并设 $y(x)=x^2v(t)$，则可化为
\[
v''(t)=1.
\]
故
\[
v(t)=\frac{t^2}{2}+C_1t+C_2,
\qquad t=\ln x.
\]
因此
\[
y(x)=x^2\left(\frac{(\ln x)^2}{2}+C_1\ln x+C_2\right).
\]
由 $y(1)=0$ 得 $C_2=0$。又
\[
y'(x)=2x\left(\frac{(\ln x)^2}{2}+C_1\ln x\right)+x(\ln x+C_1),
\]
故 $y'(1)=C_1=1$。于是
\ansmath{y(x)=x^2\left(\frac{(\ln x)^2}{2}+\ln x\right).}
\end{enumerate}
\end{solution}
\end{question}

\begin{question}
讨论函数项级数
\[
\sum_{n=1}^{\infty}\frac{x^n}{1+n},
\qquad 0\le x\le 1,
\]
在下列区间上的一致收敛性：
\begin{enumerate}[label=（\arabic*）,leftmargin=2.8em,itemsep=0.2em]
\item $[0,a]$，其中 $0<a<1$；
\item $[0,1]$.
\end{enumerate}
\begin{solution}

\begin{enumerate}[label=（\arabic*）,leftmargin=2.8em,itemsep=0.2em]
\item 当 $0\le x\le a<1$ 时，
\[
0\le \frac{x^n}{n+1}\le a^n.
\]
而 $\sum a^n$ 收敛，故由 Weierstrass 判别法，原级数在 $[0,a]$ 上一致收敛。

\item 在 $x=1$ 处，级数化为
\[
\sum_{n=1}^{\infty}\frac1{n+1},
\]
发散。因此原级数在 $[0,1]$ 上甚至不逐点收敛，更不可能一致收敛。
\end{enumerate}
故结论为
\ansmath{\text{在 }[0,a]\ (0<a<1)\text{ 上一致收敛，在 }[0,1]\text{ 上不一致收敛。}}
\end{solution}
\end{question}

\begin{question}
求函数
\[
f(x)=\frac{\ln(1+x)}{1-x}
\]
在 $x=0$ 附近的幂级数展开，并给出收敛半径。
\begin{solution}

在 $|x|<1$ 内，
\[
\ln(1+x)=\sum_{k=1}^{\infty}\frac{(-1)^{k-1}}{k}x^k,
\qquad
\frac{1}{1-x}=\sum_{m=0}^{\infty}x^m.
\]
故 Cauchy 乘积给出
\[
f(x)=\sum_{n=1}^{\infty}\left(\sum_{k=1}^{n}\frac{(-1)^{k-1}}{k}\right)x^n.
\]
前几项为
\[
f(x)=x+\frac{x^2}{2}+\frac{5x^3}{6}+\frac{7x^4}{12}+\cdots.
\]
因此
\ansmath{\frac{\ln(1+x)}{1-x}=
\sum_{n=1}^{\infty}\left(\sum_{k=1}^{n}\frac{(-1)^{k-1}}{k}\right)x^n,
\quad |x|<1.}
离原点最近的奇点为 $x=1$ 与 $x=-1$，故收敛半径
\ansmath{R=1.}
\end{solution}
\end{question}

\begin{question}
求函数 $f(x)=|\sin x|$ 在区间 $[-\pi,\pi]$ 上的 Fourier 级数，并由此求和
\[
\sum_{n=1}^{\infty}\frac{1}{4n^2-1}.
\]
\begin{solution}

$f$ 为偶函数，只含余弦项。常数项
\[
a_0=\frac{2}{\pi}\int_0^{\pi}|\sin x|\,dx=\frac{4}{\pi}.
\]
又对 $n\ge1$，利用 $|\sin x|=\sin x\ (0<x<\pi)$，得
\[
a_n=\frac{2}{\pi}\int_0^{\pi}\sin x\cos nx\,dx.
\]
由积化和差可知奇数 $n$ 时 $a_n=0$；令 $n=2m$，则
\[
a_{2m}=-\frac{4}{\pi(4m^2-1)}.
\]
故 Fourier 级数为
\ansmath{|\sin x|=\frac{2}{\pi}-\frac{4}{\pi}\sum_{m=1}^{\infty}\frac{\cos(2mx)}{4m^2-1}.}
令 $x=0$，得
\[
1=\frac{2}{\pi}-\frac{4}{\pi}\sum_{m=1}^{\infty}\frac{1}{4m^2-1}.
\]
应注意上式常数项应为 $a_0/2=2/\pi$，故整理得
\[
1=\frac{2}{\pi}+\frac{4}{\pi}\sum_{m=1}^{\infty}\frac{1}{1-4m^2}
=\frac{2}{\pi}-\frac{4}{\pi}\sum_{m=1}^{\infty}\frac{1}{4m^2-1}.
\]
从而
\ansmath{\sum_{m=1}^{\infty}\frac{1}{4m^2-1}=\frac12.}
\end{solution}
\end{question}

\begin{question}
设参数积分
\[
F(a)=\int_0^1 \frac{x^a-1}{\ln x}\,dx,
\qquad a>-1.
\]
求 $F(a)$。
\begin{solution}

对参数求导：
\[
F'(a)=\int_0^1 x^a\,dx=\frac{1}{a+1}.
\]
因此
\[
F(a)=\ln(a+1)+C.
\]
由 $F(0)=0$ 得 $C=0$，故
\ansmath{F(a)=\ln(a+1),\qquad a>-1.}
\end{solution}
\end{question}

\begin{question}
计算广义积分
\[
I=\int_0^{\infty}\frac{\arctan x}{x(1+x^2)}\,dx.
\]
\begin{solution}

令 $x=\tan t$，则 $t\in(0,\pi/2)$，并且
\[
dx=\sec^2 t\,dt,
\qquad x(1+x^2)=\tan t\sec^2 t.
\]
故
\[
I=\int_0^{\pi/2} t\cot t\,dt.
\]
分部积分，取
\[
u=t,
\qquad dw=\cot t\,dt,
\qquad w=\ln(\sin t).
\]
于是
\[
I=\bigl[t\ln(\sin t)\bigr]_0^{\pi/2}-\int_0^{\pi/2}\ln(\sin t)\,dt.
\]
边界项为 $0$，再用经典积分
\[
\int_0^{\pi/2}\ln(\sin t)\,dt=-\frac{\pi}{2}\ln 2,
\]
故
\ansmath{I=\frac{\pi}{2}\ln 2.}
\end{solution}
\end{question}


